\documentclass[11pt,a4paper]{article}

% ---------- fontes e língua ----------
\usepackage[T1]{fontenc}
\usepackage[utf8]{inputenc}
\usepackage{lmodern}
\usepackage{tgpagella}
\usepackage{tgheros}
\usepackage[portuguese]{babel}
\usepackage{microtype}

% ---------- página e cores ----------
\usepackage[a4paper,top=2.6cm,bottom=2.4cm,left=2.4cm,right=2.4cm]{geometry}
\usepackage[table]{xcolor}
\definecolor{azul}{RGB}{28,62,112}
\definecolor{azulclaro}{RGB}{226,234,246}
\definecolor{cinza}{RGB}{95,102,115}
\definecolor{laranja}{RGB}{190,105,25}
\definecolor{laranjaclaro}{RGB}{253,241,226}

% ---------- tabelas, listas, caixas ----------
\usepackage{booktabs}
\usepackage{array}
\usepackage{tabularx}
\usepackage{enumitem}
\usepackage{amssymb}
\usepackage{titlesec}
\usepackage{fancyhdr}
\usepackage{lastpage}
\usepackage{tcolorbox}

\setlength{\parindent}{0pt}
\setlength{\parskip}{6pt}

\setlist{nosep, leftmargin=*, itemsep=3pt}
\setlist[itemize]{label=\textcolor{azul}{\textbullet}}
\setlist[enumerate]{label=\textcolor{azul}{\sffamily\bfseries\arabic*.}}

\titleformat{\section}{\sffamily\Large\bfseries\color{azul}}{\thesection}{0.7em}{}[{\color{azul}\titlerule[0.6pt]}]
\titleformat{\subsection}{\sffamily\large\bfseries\color{azul}}{\thesubsection}{0.6em}{}
\titlespacing*{\section}{0pt}{1.6em}{0.8em}
\titlespacing*{\subsection}{0pt}{1.1em}{0.3em}

\pagestyle{fancy}
\setlength{\headheight}{15pt}
\fancyhf{}
\renewcommand{\headrule}{{\color{azul}\hrule width\headwidth height\headrulewidth}}
\renewcommand{\headrulewidth}{0.4pt}
\fancyhead[L]{\sffamily\small\color{cinza} Universidade da Beira Interior}
\fancyhead[R]{\sffamily\small\color{cinza} Engenharia de Software \textperiodcentered{} 2026/27}
\fancyfoot[L]{\sffamily\small\color{cinza} Guia de Jira}
\fancyfoot[R]{\sffamily\small\color{cinza} Página \thepage\ de \pageref{LastPage}}

\fancypagestyle{primeira}{%
  \fancyhf{}%
  \renewcommand{\headrulewidth}{0pt}%
  \fancyfoot[L]{\sffamily\small\color{cinza} Guia de Jira}%
  \fancyfoot[R]{\sffamily\small\color{cinza} Página \thepage\ de \pageref{LastPage}}%
}

\newcolumntype{Y}{>{\raggedright\arraybackslash}X}
\newcolumntype{L}[1]{>{\raggedright\arraybackslash}p{#1}}
\newcolumntype{N}[1]{>{\raggedright\arraybackslash\sffamily\bfseries}p{#1}}
\newcommand{\th}[1]{\textcolor{white}{\sffamily\bfseries #1}}

% nome de botão ou menu do Jira
\newcommand{\btn}[1]{{\setlength{\fboxsep}{1.6pt}\colorbox{azulclaro}{\sffamily\small\strut #1}}}
\newcommand{\seta}{\,$\rightarrow$\,}

% tabelas com cabeçalho azul e linhas alternadas
\newcommand{\tabelaini}{\begingroup\small\renewcommand{\arraystretch}{1.45}\rowcolors{2}{azulclaro!55}{white}}
\newcommand{\tabelafim}{\endgroup}

\tcbset{
  caixa/.style={colback=white, colframe=azul, coltitle=white, colbacktitle=azul,
    fonttitle=\sffamily\bfseries, arc=3pt, boxrule=0.6pt,
    left=9pt, right=9pt, top=6pt, bottom=6pt},
  dica/.style={colback=azulclaro!60, colframe=azul!70, coltitle=white, colbacktitle=azul!80,
    fonttitle=\sffamily\bfseries\small, arc=3pt, boxrule=0.5pt,
    left=9pt, right=9pt, top=5pt, bottom=5pt},
  atencao/.style={colback=laranjaclaro, colframe=laranja, coltitle=white, colbacktitle=laranja,
    fonttitle=\sffamily\bfseries\small, arc=3pt, boxrule=0.5pt,
    left=9pt, right=9pt, top=5pt, bottom=5pt}
}
\newtcolorbox{codigo}{colback=black!4, colframe=black!25, boxrule=0.4pt, arc=2pt,
  left=8pt, right=8pt, top=4pt, bottom=4pt, fontupper=\ttfamily\small}

\begin{document}
\thispagestyle{primeira}

% ---------- cabeçalho do documento ----------
\begin{tcolorbox}[colback=azul, colframe=azul, coltext=white, arc=3pt, boxrule=0pt,
    left=16pt, right=16pt, top=14pt, bottom=14pt]
  \sffamily
  {\footnotesize\bfseries\MakeUppercase{Universidade da Beira Interior}}\\[2pt]
  {\footnotesize Faculdade de Engenharia \textperiodcentered{} Departamento de Informática}\\[16pt]
  {\fontsize{28}{32}\selectfont\bfseries Guia de Jira}\\[4pt]
  {\large Do zero ao primeiro sprint \textperiodcentered{} Engenharia de Software 2026/27}
\end{tcolorbox}

Nenhum de nós usou Jira antes, e isso não é problema. Este guia explica tudo desde o princípio: o que é o Jira, como o montar, como usá-lo todos os dias e o que tem de ficar lá registado para o Scrum contar na avaliação. Não precisas de ler tudo de seguida. Faz a parte 2 uma vez (uma pessoa só) e volta à parte 4 sempre que tiveres dúvidas durante um sprint.

\begin{tcolorbox}[caixa, title={O essencial em poucas linhas}, colback=azulclaro!60]
  \begin{itemize}
    \item Cada coisa que há para fazer é um \textbf{cartão}. Se não está num cartão, para o Jira não existe.
    \item Os cartões vivem no \textbf{Backlog}. No Planning escolhemos quais entram no sprint.
    \item Durante o sprint trabalhamos no \textbf{Board}: arrastas o cartão de \textit{To Do} para \textit{In Progress} e, no fim, para \textit{Done}.
    \item Todos os cartões têm responsável, estimativa e descrição clara.
    \item Os registos do Scrum (dailies, planning, review, retrospetiva) também ficam no Jira. O enunciado pede essas provas.
    \item O que decidirmos fica escrito. O que ficar só numa conversa de WhatsApp perde-se.
  \end{itemize}
\end{tcolorbox}

%--------------------------------------------------------------------
\section{Antes de começar}

\subsection{O que é o Jira}

O Jira é uma ferramenta para organizar trabalho em equipa. Pensa num quadro de post-its muito organizado: cada post-it é um cartão com uma tarefa, alguém responsável e um estado (por fazer, a fazer, feito). O Jira junta a isso tudo o que o Scrum precisa: lista de prioridades (backlog), sprints, quadro e relatórios.

\subsection{Os nomes mudam, a ideia é a mesma}

O Jira muda os nomes das coisas de vez em quando e pode aparecer em português ou em inglês. Se um menu tiver outro nome, procura pela função e não pela palavra exata. Neste guia usamos estes nomes:

\tabelaini
\begin{tabularx}{\textwidth}{N{2.6cm} L{4.4cm} Y}
  \rowcolor{azul} \th{No guia} & \th{Pode aparecer como} & \th{O que é} \\
  Projeto & \textit{Project}, \textit{Space} (espaço) & A «casa» onde vive tudo o que é da equipa. \\
  Cartão & \textit{Issue}, \textit{Work item} & Uma coisa para fazer ou para registar. \\
  Tipo de cartão & \textit{Issue type}, \textit{Work type} & Epic, Story, Task, Bug ou Subtask. \\
  Backlog & \textit{Backlog} & A lista de tudo o que falta fazer, por ordem de prioridade. \\
  Sprint & \textit{Sprint} & Um período fechado, com objetivo, em que fazemos parte do backlog. \\
  Quadro & \textit{Board} & As colunas \textit{To Do}, \textit{In Progress} e \textit{Done} com os cartões do sprint. \\
  Responsável & \textit{Assignee} & A pessoa que vai fazer o cartão. \\
  Estimativa & \textit{Story point estimate} & Um número que mede o esforço do cartão. \\
  \bottomrule
\end{tabularx}
\tabelafim

\begin{tcolorbox}[dica, title={Dica: mudar a língua do Jira}]
  Clica na tua fotografia ou nas iniciais, no canto superior direito, e procura as definições pessoais (\textit{Personal settings} ou \textit{Account settings}). Lá há uma opção de idioma. Recomendamos inglês, porque é como o enunciado e os slides falam, mas cada um escolhe o que lhe for mais confortável.
\end{tcolorbox}

%--------------------------------------------------------------------
\section{Montar o Jira (só uma vez)}

Esta parte faz-se uma única vez e por \textbf{uma pessoa}. Sugerimos o Product Owner do Sprint 1. Os restantes só precisam de aceitar o convite no ponto 2.3.

\subsection{Criar a conta e o site}

\begin{enumerate}
  \item Abre o site do Jira (\texttt{atlassian.com/software/jira}) e carrega no botão para começar de graça (\btn{Get it free}).
  \item Regista-te com um email que não vás perder. Uma conta pessoal serve; o email da UBI também.
  \item O Jira pede um nome para o \textbf{site}. Fica com o endereço \texttt{nome.atlassian.net}. Escolhe algo simples e sem nomes de pessoas, por exemplo o nome da equipa.
  \item Escolhe o produto \textbf{Jira} e o plano \textbf{Free}. Não deves precisar de cartão de crédito. Se pedirem, pára e confirma que escolheste mesmo o plano gratuito.
  \item Guarda bem o email e a palavra-passe desta conta. É a conta «dona» do site.
\end{enumerate}

\begin{tcolorbox}[atencao, title={Atenção}]
  O plano gratuito aceita até 10 utilizadores. Somos 5 alunos e mais os docentes, por isso chega. A Atlassian pode desativar sites gratuitos que fiquem muito tempo sem atividade, por isso entrem no Jira regularmente.
\end{tcolorbox}

\subsection{Criar o projeto}

\begin{enumerate}
  \item No Jira, abre o menu \btn{Projects} (ou \btn{Spaces}) e escolhe \btn{Create project} (ou \btn{Create space}).
  \item Escolhe o modelo \btn{Scrum}.
  \item Se o Jira perguntar o tipo, escolhe \btn{Team-managed} (gerido pela equipa). É mais simples de configurar e é o que recomendamos.
  \item Dá o nome \textit{Engenharia de Software} ou outro que o grupo prefira.
  \item Escolhe a \textbf{key}: 2 a 5 letras maiúsculas. A nossa é \texttt{EDS}. É ela que dá o código aos cartões (\texttt{EDS-1}, \texttt{EDS-2}, \dots) e que depois usamos nas branches do GitHub. Pensa bem, porque é difícil mudar mais tarde.
  \item No acesso, deixa o projeto \textbf{Private}, para só entrar quem convidarmos.
  \item Cria. O Jira abre no Backlog.
\end{enumerate}

\begin{tcolorbox}[dica, title={Dica: se não vês Backlog ou Sprints}]
  Abre as definições do projeto (\btn{Space settings} ou \btn{Project settings}, na barra lateral) e procura \btn{Features}. Liga \btn{Backlog} e \btn{Sprints}. O modelo Scrum costuma trazer isto ligado.
\end{tcolorbox}

\subsection{Convidar a equipa e os docentes}

O enunciado exige que os docentes tenham acesso ao Jira, por isso este passo conta para a nota.

\begin{enumerate}
  \item Na barra lateral, abre as definições do projeto e escolhe \btn{Access}.
  \item Carrega em \btn{Add people}, escreve o email de cada pessoa e escolhe a função.
  \item Cada pessoa recebe um email com o convite. Tem de o aceitar e, se não tiver conta Atlassian, criar uma. Se o email não chegar, vê o spam.
  \item No fim, cada um confirma que consegue abrir o Backlog e o Board.
\end{enumerate}

\tabelaini
\begin{tabularx}{\textwidth}{N{2.8cm} Y L{4.6cm}}
  \rowcolor{azul} \th{Função} & \th{O que pode fazer} & \th{Quem} \\
  Administrator & Mexer nas definições, nos acessos e até apagar o projeto. & Só duas pessoas: o PO e o SM do Sprint 1. \\
  Member & Criar, editar, comentar e mover cartões. & Toda a equipa. \\
  Viewer & Ver o projeto e os cartões, e pouco mais. & Os docentes. Se quiserem comentar, passam a Member. \\
  \bottomrule
\end{tabularx}
\tabelafim

Pede aos docentes (Nuno Pombo e Carlos Cunha) o email que usam para o convite.

\begin{tcolorbox}[atencao, title={Atenção}]
  Quanto menos administradores, melhor: quem é Administrator pode apagar o projeto inteiro. Duas pessoas chegam e assim há sempre uma de reserva.
\end{tcolorbox}

\subsection{Montar a estrutura que o enunciado pede}

O enunciado sugere sete «quadros»: \textit{Product Backlog}, \textit{To Do}, \textit{In Progress}, \textit{Done}, \textit{Artifacts}, \textit{Daily SCRUM} e \textit{Team Contract}. Os primeiros quatro já existem no Jira. Para os outros três usamos \textbf{epics}, que funcionam como pastas.

\tabelaini
\begin{tabularx}{\textwidth}{L{4.2cm} Y}
  \rowcolor{azul} \th{O enunciado pede} & \th{Como fazemos no Jira} \\
  Product Backlog & O ecrã \btn{Backlog}. Já existe. \\
  To Do, In Progress, Done & As três colunas do \btn{Board}. Já existem. \\
  Team Contract & Um epic chamado \textit{Team Contract}. \\
  Artifacts & Um epic chamado \textit{Artifacts}. \\
  Daily SCRUM & Um epic chamado \textit{Daily SCRUM}. \\
  \bottomrule
\end{tabularx}
\tabelafim

Para criar cada epic:

\begin{enumerate}
  \item Carrega em \btn{Create} (botão azul no topo).
  \item Em tipo de cartão, escolhe \btn{Epic}.
  \item Escreve o nome exatamente como na tabela e carrega em \btn{Create}.
\end{enumerate}

O que pomos em cada um:

\begin{itemize}
  \item \textbf{Team Contract:} na descrição, uma linha a dizer desde quando está em vigor. Anexa o PDF do contrato e o do Registo de Uso de IA (abre o epic e arrasta os ficheiros, ou usa o clipe, \btn{Attach}). Sempre que o contrato mudar, anexamos a nova versão e comentamos o que mudou e porquê.
  \item \textbf{Artifacts:} ligações para o repositório do GitHub, os wireframes, os mockups e os diagramas. Podes colar os links na descrição. Para cada artefacto grande, cria um cartão filho (por exemplo \textit{Wireframes v1}) com o link ou o ficheiro.
  \item \textbf{Daily SCRUM:} um cartão filho por sprint para os dailies (\textit{Daily – Sprint 1}) e um para cada cerimónia (\textit{Planning – Sprint 1}, \textit{Review – Sprint 1}, \textit{Retrospetiva – Sprint 1}). Cada um fica dentro do sprint a que pertence e passa a \textit{Feito} quando a reunião acontece e fica registada.
\end{itemize}

Para criar um cartão filho dentro de um epic, abre o epic e usa o botão para adicionar um cartão filho (\textit{child}). Se não o encontrares, cria o cartão normalmente e escolhe o epic no campo \btn{Parent}.

\begin{tcolorbox}[dica, title={Dica: links úteis na barra lateral}]
  Se na barra lateral aparecer \btn{Add shortcut} ou \btn{Add link}, usa-o para pôr lá o link do repositório do GitHub. Fica sempre à mão. O Jira criou também um espaço no Confluence com o mesmo nome do projeto (\texttt{EDS}). Tem páginas para o contrato, a equipa, o SCRUM e os artefactos. As atas ficam nos cartões e as páginas servem de resumo.
\end{tcolorbox}

\subsection{Ligar o Jira ao GitHub (recomendado)}

Quando a ligação está feita, o Jira mostra dentro de cada cartão as branches, os commits e os \textit{pull requests} que lhe dizem respeito. Isto ajuda a mostrar a rastreabilidade entre o que pedimos e o que foi feito.

\begin{enumerate}
  \item No Jira, abre as aplicações (\btn{Apps}, ou \btn{Explore more apps}) e procura \textbf{GitHub for Jira}.
  \item Instala a aplicação (\btn{Get app}) e segue o assistente: \btn{Connect GitHub organization}, entra no GitHub e escolhe o repositório da equipa.
  \item Quem faz esta ligação tem de ter permissões no repositório, por isso é quem o criou.
\end{enumerate}

Se o assistente der problemas, não bloqueia nada: o que importa é seguirmos a regra do ponto 5 (escrever a key do cartão nos nomes). A ligação só torna isso visível no Jira.

\subsection{Montar o painel (opcional)}

O enunciado fala em criar um «dashboard» com as secções \textit{Product Backlog}, \textit{To Do}, \textit{In Progress}, \textit{Done}, \textit{Artifacts}, \textit{Daily SCRUM} e \textit{Team Contract}. Uma forma de o ter num só ecrã é guardar um filtro por secção e pôr cada um num painel.

\begin{enumerate}
  \item No Jira, abre a pesquisa avançada (\btn{Filters}), cola a pesquisa da tabela, carrega em \btn{Save as} e dá o nome da secção.
  \item Depois abre \btn{Dashboards}, cria um painel novo e acrescenta, para cada filtro, o gadget \btn{Filter Results}.
\end{enumerate}

\tabelaini
\begin{tabularx}{\textwidth}{N{3.2cm} Y}
  \rowcolor{azul} \th{Secção} & \th{Pesquisa} \\
  Product Backlog & \small\texttt{project = EDS AND issuetype != Epic AND sprint is EMPTY AND statusCategory != Done ORDER BY Rank} \\
  To Do & \small\texttt{project = EDS AND sprint in openSprints() AND statusCategory = 'To Do'} \\
  In Progress & \small\texttt{project = EDS AND sprint in openSprints() AND statusCategory = 'In Progress'} \\
  Done & \small\texttt{project = EDS AND sprint in openSprints() AND statusCategory = Done} \\
  Artifacts & \small\texttt{project = EDS AND parent = EDS-2} \\
  Daily SCRUM & \small\texttt{project = EDS AND parent = EDS-3} \\
  Team Contract & \small\texttt{project = EDS AND parent = EDS-1} \\
  \bottomrule
\end{tabularx}
\tabelafim

%--------------------------------------------------------------------
\section{Conhecer o ecrã}

\subsection{A barra lateral}

Quando entras no projeto, a barra do lado esquerdo tem os sítios principais. Os nomes podem variar um pouco.

\tabelaini
\begin{tabularx}{\textwidth}{N{3cm} Y}
  \rowcolor{azul} \th{Sítio} & \th{Para que serve} \\
  Summary & Resumo do projeto e da atividade recente. \\
  Timeline & Os epics numa linha do tempo, tipo roadmap. \\
  Backlog & A lista de tudo o que falta fazer e a preparação dos sprints. É o ecrã do Product Owner. \\
  Board & O quadro do sprint em curso, com as colunas. É onde passamos o dia a dia. \\
  Reports & Relatórios do sprint (por exemplo o \textit{burndown}). Servem de prova na Review. \\
  Settings & Definições do projeto, acessos e funcionalidades. \\
  \bottomrule
\end{tabularx}
\tabelafim

\begin{tcolorbox}[dica, title={Dica: atalhos de teclado}]
  Dentro do Jira, a tecla \texttt{c} costuma abrir a criação de um cartão e a tecla \texttt{/} abre a pesquisa. A tecla \texttt{[} esconde e mostra a barra lateral.
\end{tcolorbox}

\subsection{Como é um cartão}

Quando abres um cartão, vês estes campos:

\tabelaini
\begin{tabularx}{\textwidth}{N{3.4cm} Y L{3.4cm}}
  \rowcolor{azul} \th{Campo} & \th{O que é} & \th{Quem preenche} \\
  Key (ex.: EDS-12) & O código único do cartão. Usamo-lo no GitHub. & O Jira \\
  Summary & O título, curto e claro. & Quem cria \\
  Description & O que é para fazer, a user story e os critérios de aceitação. & Quem cria, com o PO \\
  Status & Em que coluna está: \textit{To Do}, \textit{In Progress} ou \textit{Done}. & Quem trabalha nele \\
  Assignee & Quem vai fazer o cartão. Uma só pessoa. & No Planning \\
  Story point estimate & A estimativa do esforço. & A equipa, em conjunto \\
  Parent / Epic & O epic a que pertence. & Quem cria \\
  Sprint & O sprint em que está. & No Planning \\
  Labels & Etiquetas livres, por exemplo \textit{frontend} ou \textit{backend}. & Quem cria \\
  Comments & A conversa sobre o cartão. & Todos \\
  Attachments & Ficheiros anexados (capturas, PDF, imagens). & Todos \\
  \bottomrule
\end{tabularx}
\tabelafim

\subsection{Os tipos de cartão}

\tabelaini
\begin{tabularx}{\textwidth}{N{2.4cm} Y L{4.8cm}}
  \rowcolor{azul} \th{Tipo} & \th{Quando usar} & \th{Exemplo} \\
  Epic & Um tema grande, que agrupa muitos cartões. & \textit{Gestão de utilizadores} \\
  Story & Uma user story, vista do ponto de vista de quem usa. & \textit{Recuperar a palavra-passe} \\
  Task & Trabalho técnico ou de organização que não é uma user story. & \textit{Configurar a integração contínua} \\
  Bug & Algo que está errado e tem de ser corrigido. & \textit{O botão «Guardar» não responde} \\
  Subtask & Um pedaço pequeno de outro cartão. & \textit{Escrever os testes do login} \\
  \bottomrule
\end{tabularx}
\tabelafim

\subsection{Os estados}

O nosso Jira está em português, por isso as colunas chamam-se \textit{A fazer}, \textit{Fazendo} e \textit{Feito}, e os tipos de cartão \textit{Tarefa} e \textit{História}. É o mesmo que \textit{To Do}, \textit{In Progress}, \textit{Done}, \textit{Task} e \textit{Story}. O ciclo de vida de um cartão tem três passos:

\begin{itemize}
  \item \textbf{To Do:} está no sprint mas ainda ninguém começou.
  \item \textbf{In Progress:} alguém está a trabalhar nele. Fica aqui até o \textit{pull request} ser aprovado e integrado.
  \item \textbf{Done:} cumpre a nossa Definição de Concluído e o Product Owner aceitou.
\end{itemize}

%--------------------------------------------------------------------
\section{Um sprint, passo a passo}

\subsection{Construir o Backlog: criar cartões}

\begin{enumerate}
  \item Carrega em \btn{Create}, no topo, ou em \btn{+ Create} no fundo da lista do Backlog.
  \item Escolhe o tipo (normalmente \btn{Story} ou \btn{Task}).
  \item Escreve o \btn{Summary}: curto, no formato «verbo + coisa». Por exemplo, \textit{Recuperar palavra-passe}.
  \item Preenche a \btn{Description} (modelo abaixo).
  \item Se couber num epic, escolhe-o em \btn{Parent}.
  \item Deixa o \btn{Assignee} vazio até ao Planning, quando se decide quem faz o quê. Carrega em \btn{Create}.
\end{enumerate}

Modelo de descrição de uma user story (o exemplo é inventado, só para mostrar o formato):

\begin{codigo}
Como utilizador registado, quero recuperar a minha palavra-passe, para voltar a entrar quando me esqueço dela.\\[4pt]
Critérios de aceitação:\\
- Há uma ligação 'Esqueci-me da palavra-passe' no ecrã de entrada.\\
- Recebo um email com uma ligação que expira ao fim de uma hora.\\
- Depois de a mudar, consigo entrar com a nova palavra-passe.
\end{codigo}

\begin{tcolorbox}[dica, title={Dica: caixas de verificação}]
  No editor de descrição do Jira, escreve \texttt{[]} seguido de um espaço para criar uma lista com caixas de seleção. É útil para os critérios de aceitação e para a Definição de Concluído.
\end{tcolorbox}

\subsection{Estimar com story points}

Os story points medem o esforço relativo de um cartão, não horas. Um cartão de 4 pontos dá cerca do dobro do trabalho de um de 2.

\begin{itemize}
  \item Usamos a escala \textbf{1, 2, 3, 5, 8, 13}.
  \item Escolhemos um cartão pequeno e consensual como referência (por exemplo, o de 2 pontos) e comparamos os outros com ele.
  \item Fazemos \textit{planning poker} simples: cada um pensa num número e todos o mostram ao mesmo tempo (dedos, ou mensagem no grupo). Se houver muita diferença, quem escolheu o mais baixo e o mais alto explica porquê, e repetimos.
  \item Um cartão de mais de 13 pontos é grande demais. Parte-se em cartões mais pequenos.
\end{itemize}

Para registar, abre o cartão e preenche \btn{Story point estimate}, ou clica na bolinha cinzenta com o número, ao lado do cartão, na lista do Backlog.

\subsection{Priorizar (o trabalho do Product Owner)}

No Backlog, o que está no topo é o que vem primeiro. O Product Owner arrasta os cartões para os pôr por ordem de valor: o mais importante em cima. Só o Product Owner mexe na ordem; os outros sugerem e ele decide, depois de ouvir a equipa.

\subsection{Sprint Planning no Jira}

Este é o primeiro dia do sprint. O Scrum Master conduz a reunião e mexe no ecrã.

\begin{enumerate}
  \item Abre o \btn{Backlog}. Se ainda não houver um sprint para este período, carrega em \btn{Create sprint}.
  \item Muda o nome para \textit{Sprint 1}, \textit{Sprint 2}, etc.
  \item O Product Owner apresenta os cartões de cima do backlog. A equipa discute e decide até onde cabe no sprint.
  \item Arrasta esses cartões do backlog para o sprint. Vê o total de story points no topo do sprint.
  \item Em cada cartão do sprint, define o \btn{Assignee}.
  \item Carrega em \btn{Start sprint}. Confirma o nome, ajusta as \textbf{datas} para as oficiais (a tabela está no contrato) e escreve o \btn{Sprint goal}, o objetivo do sprint numa frase. Termina com «PO: nome, SM: nome», para ficar registado quem tem cada papel.
  \item Escreve a ata curta no cartão \textit{Planning – Sprint N} (epic \textit{Daily SCRUM}): data, quem esteve, objetivo e cartões escolhidos. Junta uma captura do ecrã do sprint.
\end{enumerate}

\begin{tcolorbox}[atencao, title={Atenção: as datas}]
  Ao iniciar o sprint, o Jira preenche as datas com o momento em que clicas. Muda-as para as datas oficiais, senão os relatórios ficam errados.
\end{tcolorbox}

\begin{tcolorbox}[dica, title={E se já estamos a meio do sprint?}]
  Começa o sprint com as datas reais e põe lá os cartões do que está em curso. Para o trabalho que já foi feito (por exemplo o contrato), cria o cartão e passa-o para \textit{Done} com um comentário a dizer o que foi feito. Não inventamos registos de dias que já passaram: só escrevemos o que realmente aconteceu.
\end{tcolorbox}

\subsection{O dia a dia no Board}

Abre o \btn{Board}. Vês as colunas e, em cada uma, os cartões do sprint.

\begin{itemize}
  \item \textbf{Pegar num cartão:} confirma que és tu o responsável, arrasta-o para \textit{In Progress} e começa.
  \item \textbf{Acabar:} depois de o \textit{pull request} estar aprovado e integrado e de o Product Owner aceitar, arrasta-o para \textit{Done}.
  \item \textbf{Atualizar no momento.} Mover o cartão leva dois segundos e evita confusões na daily.
  \item \textbf{No máximo dois cartões em \textit{In Progress}} por pessoa.
  \item \textbf{Bloqueado?} Abre o cartão, escreve um comentário a explicar o bloqueio e menciona quem pode ajudar com \texttt{@nome}. Podes ainda marcá-lo com \btn{Add flag} (no menu de três pontos do cartão): fica destacado no quadro.
  \item \textbf{Dividir o trabalho:} se um cartão tem partes distintas, cria subtasks e atribui cada uma a quem a vai fazer.
  \item \textbf{Mostrar o trabalho:} anexa capturas de ecrã e links nos comentários do cartão. Quanto mais claro, mais fácil é aceitar na Review.
\end{itemize}

\subsection{O Daily Scrum e o registo no Jira}

Fazemos a daily como o contrato diz: 15 minutos, em pé, com o Board aberto. O Scrum Master vai passando por cada pessoa e pelos cartões. A seguir, \textbf{cada um escreve o seu registo} no cartão \textit{Daily – Sprint N}, como comentário. Cola este modelo:

\begin{codigo}
Dia 14/10 - Nome\\
Feito: o que fiz desde a última daily\\
A fazer: o que vou fazer hoje\\
Bloqueios: o que me está a impedir (ou 'nenhum')
\end{codigo}

Quem não pôde estar escreve na mesma, no mesmo dia. O Scrum Master confirma, até ao fim do dia, que está tudo escrito.

\subsection{Refinement}

A meio da semana, o Product Owner e a equipa abrem o Backlog e olham para os cartões que se seguem. Completam as descrições, acrescentam critérios de aceitação, estimam o que falta e partem os cartões grandes. Registar uma nota curta em \textit{Planning – Sprint N} basta (o que se afinou, o que ficou por decidir).

\subsection{Review, Retrospetiva e fechar o sprint}

\begin{enumerate}
  \item \textbf{Antes da Review:} no Board, confirma que cada cartão em \textit{Done} cumpre mesmo a Definição de Concluído.
  \item \textbf{Durante a Review:} mostramos o que ficou pronto. O Product Owner aceita ou não os cartões. Quem não foi aceite volta para \textit{To Do} ou \textit{In Progress}.
  \item \textbf{Provas:} tira capturas do Board no fim do sprint e dos relatórios (\btn{Reports}, por exemplo o \textit{Sprint report} ou o \textit{Burndown chart}, se existirem). Anexa-as ao cartão \textit{Review – Sprint N}.
  \item \textbf{Retrospetiva:} no cartão \textit{Retrospetiva – Sprint N}, escrevemos o que correu bem, o que correu mal e duas ou três melhorias para o sprint seguinte. Cada melhoria tem um responsável.
  \item \textbf{Fechar o sprint:} no Board ou no Backlog, carrega em \btn{Complete sprint}. O Jira pergunta o que fazer aos cartões que não ficaram \textit{Done}. Escolhe \btn{Backlog}: como o contrato diz, o que não ficou pronto volta ao Product Backlog para ser priorizado de novo, e o sprint não se prolonga.
\end{enumerate}

\subsection{O que tem de ficar no Jira em cada sprint}

O enunciado avalia o Scrum através do que fica registado, por isso cada sprint deve ter:

\begin{itemize}[label=\textcolor{azul}{$\square$}]
  \item Sprint com nome, datas oficiais e objetivo, com o PO e o SM indicados.
  \item Cartões do sprint com responsável, estimativa e descrição.
  \item Ata do Planning, com captura do ecrã.
  \item Um registo diário de cada pessoa, em cada dia útil, no cartão \textit{Daily – Sprint N}.
  \item Nota do Refinement.
  \item Review, com captura do Board e dos relatórios.
  \item Retrospetiva, com as melhorias e os responsáveis.
  \item Backlog atualizado, com o que ficou por fazer.
\end{itemize}

%--------------------------------------------------------------------
\section{Jira e GitHub juntos}

Cada cartão tem um código (\textit{key}), por exemplo \texttt{EDS-12}. Se o escrevermos nos sítios certos, o Jira liga o trabalho ao cartão sozinho.

\begin{itemize}
  \item \textbf{Branch:} começa pela key, em maiúsculas, dentro do nome da branch: \texttt{feature/EDS-12-login}.
  \item \textbf{Commit:} começa com a key, seguida de uma frase curta: \texttt{EDS-12 faz o ecrã de entrada}.
  \item \textbf{Pull request:} a key vai no título, e a descrição diz o que foi feito.
\end{itemize}

\begin{codigo}
git switch dev\\
git pull\\
git switch -c feature/EDS-12-login
\end{codigo}

A key tem de estar em maiúsculas (\texttt{EDS-12}, e não \texttt{eds-12}). Depois de fazeres \textit{push}, a ligação pode demorar alguns minutos a aparecer no cartão. Nenhum pull request é aprovado por quem o escreveu, como o contrato diz.

%--------------------------------------------------------------------
\section{Quem faz o quê}

\tabelaini
\begin{tabularx}{\textwidth}{N{3.2cm} Y}
  \rowcolor{azul} \th{Papel} & \th{No Jira} \\
  Product Owner & Mantém o Backlog por ordem de prioridade. Escreve ou revê as user stories e os critérios de aceitação. Aceita (ou recusa) os cartões na Review. \\
  Scrum Master & Cria o sprint e carrega em \btn{Start sprint} e \btn{Complete sprint}. Confirma que as dailies estão escritas. Escreve as atas e junta as capturas. Medeia se houver dúvidas sobre o uso do Jira. \\
  Toda a equipa & Cria cartões, estima, escolhe e move os seus cartões, comenta, escreve a sua daily, mantém tudo atualizado e põe a key do cartão no GitHub. \\
  Administradores & Convidam e retiram pessoas e tratam das definições. Mais nada. \\
  \bottomrule
\end{tabularx}
\tabelafim

%--------------------------------------------------------------------
\section{Problemas comuns}

\tabelaini
\begin{tabularx}{\textwidth}{L{4.6cm} Y}
  \rowcolor{azul} \th{O que se passa} & \th{O que fazer} \\
  O Board está vazio & O sprint ainda não foi iniciado, ou os cartões estão no backlog e não no sprint. Vai ao Backlog, arrasta os cartões para o sprint e carrega em \btn{Start sprint}. \\
  Não consigo arrastar nem editar nada & Provavelmente tens a função \textit{Viewer}. Pede a um administrador para te passar a \textit{Member}. \\
  Um cartão desapareceu & Pode estar noutro sprint, noutro epic, ou escondido por um filtro. Usa a pesquisa (tecla \texttt{/}) com o nome ou a key. \\
  Criei um cartão no sítio errado & Abre-o e corrige os campos \btn{Sprint} e \btn{Parent}, ou arrasta-o para o sítio certo no Backlog. \\
  Apaguei um cartão sem querer & Não dá para desfazer. Por isso não apagamos cartões: se for um engano, avisa o Scrum Master. \\
  O convite não chegou & Procura no spam. Se não estiver lá, o administrador reenvia-o. \\
  Recebo emails a mais & Nas definições pessoais da tua conta há opções de notificações. Podes reduzi-las, mas não ignores as menções com \texttt{@}. \\
  A ligação ao GitHub não aparece & Confirma se a key está em maiúsculas, se fizeste \textit{push} e se esperaste uns minutos. \\
  Não encontro um botão & Procura pelo que ele faz, não pelo nome. Os nomes mudam. Se continuar a faltar, pergunta no grupo. \\
  \bottomrule
\end{tabularx}
\tabelafim

%--------------------------------------------------------------------
\section{Checklist de arranque}

\subsection{A pessoa que cria o Jira}

\begin{itemize}[label=\textcolor{azul}{$\square$}]
  \item Criar a conta Atlassian e o site (plano \textit{Free}).
  \item Criar o projeto com o modelo Scrum, \textit{team-managed}, e escolher a key.
  \item Convidar a equipa (\textit{Member}) e os docentes (\textit{Viewer}).
  \item Dar \textit{Administrator} a mais uma pessoa.
  \item Criar os epics: \textit{Team Contract}, \textit{Artifacts}, \textit{Daily SCRUM}, \textit{Presenças e faltas} e \textit{Equipa e papéis}.
  \item Anexar o contrato e o registo de uso de IA ao epic \textit{Team Contract}.
  \item Ligar ao GitHub (opcional) e pôr o link do repositório no epic \textit{Artifacts}.
\end{itemize}

\subsection{Cada elemento da equipa}

\begin{itemize}[label=\textcolor{azul}{$\square$}]
  \item Aceitar o convite e criar a conta Atlassian.
  \item Confirmar que consegue abrir o Backlog e o Board.
  \item Mudar a língua para a que lhe for mais confortável.
  \item Treinar num cartão de teste: o Scrum Master cria um, cada um comenta nele e arrasta-o no Board, e no fim o Scrum Master apaga-o.
  \item Ler as partes 4 e 5 antes do primeiro Planning.
\end{itemize}

\subsection{Para o primeiro Planning}

\begin{itemize}[label=\textcolor{azul}{$\square$}]
  \item Backlog com as primeiras user stories, com descrição e critérios de aceitação.
  \item Cartões estimados em story points.
  \item Backlog ordenado pelo Product Owner.
  \item Sprint criado, com datas oficiais, objetivo, PO e SM.
  \item Cartões do sprint com responsável.
  \item Cartões \textit{Planning}, \textit{Daily}, \textit{Review} e \textit{Retrospetiva} criados dentro do epic \textit{Daily SCRUM} e do sprint.
\end{itemize}

%--------------------------------------------------------------------
\section{Glossário}

\tabelaini
\begin{tabularx}{\textwidth}{N{3.4cm} Y}
  \rowcolor{azul} \th{Termo} & \th{Significado} \\
  Assignee & Responsável por um cartão. \\
  Backlog & Lista de tudo o que falta fazer, por ordem de prioridade. \\
  Board & Quadro com as colunas do sprint em curso. \\
  Burndown chart & Gráfico que mostra o trabalho que falta ao longo do sprint. \\
  Epic & Cartão grande que agrupa outros cartões. \\
  Flag & Marca que destaca um cartão bloqueado. \\
  Key & Código único de um cartão, por exemplo EDS-12. \\
  Label & Etiqueta livre para organizar cartões. \\
  Sprint & Período fechado, com objetivo, para fazer parte do backlog. \\
  Sprint goal & O objetivo do sprint, numa frase. \\
  Story & Cartão que descreve uma necessidade do ponto de vista de quem usa. \\
  Story points & Medida relativa do esforço de um cartão. \\
  Subtask & Parte pequena de outro cartão. \\
  Velocity & Quantidade de story points que a equipa costuma fechar num sprint. \\
  Work item & O novo nome de \textit{issue}: o mesmo que cartão. \\
  \bottomrule
\end{tabularx}
\tabelafim

\end{document}
